Probe selection precedes primary scanning and selects inverted lists,
not database-vector survivors. For IVF centroid $\mu_i$, the selector
orders the FP32-valued computed routing scores
\begin{equation}
 s_i=\|\mu_i\|_2^2-2\langle q',\mu_i\rangle,
 \qquad (s_i,i)\ \text{ordered lexicographically}.
 \label{eq:5-probe-score}
\end{equation}
Omitting the common query norm preserves centroid ordering. In contrast
to primary squared-distance tables, these scores can be negative. The output is
\[
 n_{\rm sel}=\min(n_{\rm probe},n_{\rm list}),\qquad
 I=\operatorname{Top}^{(s_i,i)}_{n_{\rm sel}}\{0,\ldots,n_{\rm list}-1\},
\]
with list-length prefix offsets.

OpenAI v69 uses TF32 rotation/centroid GEMMs with FP32 accumulation. Probe exactness concerns their computed scores; FP32 \texttt{fma} retains earlier rounding (\cref{sec:7-precision}).

\subsection{Radix selection followed by deterministic ordering}

One block makes four MSB-first radix passes. For prefix-restricted histogram $H_t$ and one-based residual rank $r_t$ ($r_0=n_{\rm sel}$),
\[
 b_t=\min\Bigl\{b:\sum_{a\leq b}H_t[a]\geq r_t\Bigr\},\qquad
 r_{t+1}=r_t-\sum_{a<b_t}H_t[a].
\]
The chosen bin extends the prefix. Four bytes fix the cutoff/equal-key count.

Canonicalise signed zeros to positive zero, then map bits $u$ by
\begin{equation}
 F(u)=\begin{cases}
       \mathop{\sim}u,&\text{if the sign bit is set},\\
       u\mathbin{|}2^{31},&\text{otherwise}.
      \end{cases}
 \label{eq:5-probe-key}
\end{equation}
For non-NaN scores, monotonicity follows by cases:
\[
 \begin{aligned}
 D_1<D_2<0&\Rightarrow u_1>u_2\Rightarrow{\sim}u_1<{\sim}u_2,\\
 0\leq D_1<D_2&\Rightarrow u_1<u_2\Rightarrow
 (u_1\mathbin{|}2^{31})<(u_2\mathbin{|}2^{31}),\\
 D_1<0\leq D_2&\Rightarrow F(u_1)<2^{31}\leq F(u_2).
 \end{aligned}
\]
Infinities/zeros are covered, unlike the nonnegative-only \cref{eq:5-packed-key}. Repeated \texttt{fma} fixes recomputed keys.

Histogram storage is
\[
 256\ \text{bins}\times32\ \text{lane replicas}\times4\ \text{bytes}=32\,\mathrm{KiB}.
\]
Replication reduces same-warp collisions, leaving inter-warp contention/reduction cost.

Let $h$ be the cutoff key. Collection forms
\[
 A_< =\{i:F(u_i)<h\},\quad A_= =\{i:F(u_i)=h\},\quad
 m=n_{\rm sel}-|A_<|.
\]
Take $A_<$ plus the $m$ smallest IDs in $A_=$, using ordered ballots/prefix counts rather than atomic arrival. For the resulting set $A$,
\[
 \operatorname{rank}(i)=|\{j\in A:(s_j,j)\prec(s_i,i)\}|,
 \qquad\mathrm{out}[\operatorname{rank}(i)]=i.
\]

Selected lengths produce prefix offsets. For excess probes,
\[
 i_j=i_{n_{\rm sel}-1},\quad L_j=0\quad(j\geq n_{\rm sel}),
 \qquad\omega_{j+1}=\omega_j+L_j.
\]
Zero lengths prevent rescans. Require $n_{\rm list},n_{\rm probe}>0$.

\subsection{Costs and separate correctness contracts}

Histogram/collection reread scores; ranking is $O(n_{\rm sel}^2)$ with growing shared memory. Small selections still pay fixed histogram storage; this is deterministic integration, not uniform linear cost.

Tests cover ties/zeros/infinities/extremes/excess probes (\cref{sec:7-4}).

FP16 approximation and cross-process checks are separate (\cref{sec:7-4}).